\SourceLeaf{228}%
\[
\cos(1)L=X,\qquad\cos(2)L=Y,\qquad\cos(3)L=Z
\]
\noindent und bezeichnen die Koordinaten des Punktes \(A\) mit \(x,y,z\). Ferner seien \(x+\mathrm{d}x,y+\mathrm{d}y,z+\mathrm{d}z\) die Koordinaten eines anderen Punktes \(A'\) auf der gekrümmten Fläche, \(\mathrm{d}s\) sein unendlich kleiner Abstand von \(A\) und \(\lambda\) der Punkt der Kugeloberfläche, der die Richtung des Elements \(AA'\) darstellt. Dann gilt
\[
\mathrm{d}x=\mathrm{d}s\mathord{.}\cos(1)\lambda,\quad\mathrm{d}y=\mathrm{d}s\mathord{.}\cos(2)\lambda,\quad\mathrm{d}z=\mathrm{d}s\mathord{.}\cos(3)\lambda
\]
\noindent und da \(\lambda L=90^0\) gelten muss,
\[
X\cos(1)\lambda+Y\cos(2)\lambda+Z\cos(3)\lambda=0
\]
\noindent Durch Kombination dieser Gleichungen erhalten wir
\[
X\mathrm{d}x+Y\mathrm{d}y+Z\mathrm{d}z=0
\]
\par\indent Es gibt zwei allgemeine Methoden, eine gekrümmte Fläche zu beschreiben. Die \textit{erste} verwendet eine Gleichung zwischen den Koordinaten \(x,y,z\), und wir nehmen an, dass sie auf die Form \(W=0\) gebracht worden ist. Dabei ist \(W\) eine Funktion der unbestimmten Größen \(x,y,z\). Das totale Differential der Funktion \(W\) sei
\[
\mathrm{d}W=P\mathrm{d}x+Q\mathrm{d}y+R\mathrm{d}z
\]
\noindent dann gilt auf der gekrümmten Fläche
\[
P\mathrm{d}x+Q\mathrm{d}y+R\mathrm{d}z=0
\]
\noindent und folglich
\[
P\cos(1)\lambda+Q\cos(2)\lambda+R\cos(3)\lambda=0
\]
\noindent Da diese Gleichung ebenso wie die oben aufgestellte für die Richtungen aller Elemente \(\mathrm{d}s\) auf der gekrümmten Fläche gelten muss, erkennt man leicht, dass \(X,Y,Z\) zu \(P,Q,R\) proportional sein müssen. Da nun
\[
XX+YY+ZZ=1
\]
\noindent gilt, haben wir entweder
\begingroup\small
\[
X=\frac{P}{\BareRoot{(PP+QQ+RR)}},\quad Y=\frac{Q}{\BareRoot{(PP+QQ+RR)}},\quad Z=\frac{R}{\BareRoot{(PP+QQ+RR)}}
\]\endgroup

\noindent oder
\begingroup\small
\[
X=\frac{-P}{\BareRoot{(PP+QQ+RR)}},\quad Y=\frac{-Q}{\BareRoot{(PP+QQ+RR)}},\quad Z=\frac{-R}{\BareRoot{(PP+QQ+RR)}}
\]\endgroup

